\section{The trace coefficient and the passage to all radii}
\label{sec:trace-asymptotic}
Two different normalizations enter the trace asymptotic. Translation Ward
identities fix the amplitude of the stress tensor, while the running Wilson
coefficient fixes the logarithmic scale. After those normalizations, the
numerical coefficient comes from a free curvature contraction. Finally one
must remove the rounded positions and the chosen sequence of scales. We
separate these operations because none can be replaced by choosing a
convenient normalization of the trace after the fact.

\subsection{The free scalar contraction}
Put $N=\dim_{\mathbb R}\mathfrak g$ and use an orthonormal color basis.
The reference potential covariance in four Euclidean dimensions is
$\delta^{AB}\delta_{ab}D(x)$, where
\[
 D(x)=\frac1{4\pi^2|x|^2},\qquad x\neq0.
\]
Indices $a,b,i,j$ range over the four coordinate directions and $A,B$
over the $N$ colors. The curvature covariance obtained by differentiating
this potential is $\delta^{AB}C_{abij}(x)$. Explicitly,
\begin{align*}
 C_{abij}(x)={}&-\delta_{bj}\partial_a\partial_iD(x)
       +\delta_{bi}\partial_a\partial_jD(x)\\
       &+\delta_{aj}\partial_b\partial_iD(x)
       -\delta_{ai}\partial_b\partial_jD(x).
\end{align*}
Write $r=|x|$ and $H_{ai}=\delta_{ai}-2x_ax_i/r^2$. Differentiating
$D$ gives
\[
 \partial_a\partial_iD(x)
   =\frac1{4\pi^2}(-2\delta_{ai}r^{-4}+8x_ax_ir^{-6}),
 \qquad
 C_{abij}(x)=\frac{H_{ai}H_{bj}-H_{aj}H_{bi}}{\pi^2r^4}.
\]
For the second identity expand the products of $H$. Their quartic terms
cancel, leaving exactly the four second-derivative terms above.
Moreover $H=I-2uu^{\mathsf T}$ for the unit vector $u=x/r$, hence
$HH^{\mathsf T}=I$. It follows that
\begin{align*}
 \sum_{i,j}(H_{ai}H_{bj}-H_{aj}H_{bi})^2
 &=2\Bigl(\sum_iH_{ai}^2\Bigr)
       \Bigl(\sum_jH_{bj}^2\Bigr)
       -2\Bigl(\sum_iH_{ai}H_{bi}\Bigr)^2\\
 &=2-2\delta_{ab}.
\end{align*}
Summing over the sixteen ordered pairs $(a,b)$ gives $24$.

The reference scalar is
\[
 O(x)=\frac14\sum_{A,a,b}:F_{ab}^A(x)F_{ab}^A(x):,
\]
where the sum is over ordered coordinate pairs. Here the colons mean the
quadratic Wick field: in its two-point function only contractions between
the two different positions are included. There are exactly two such
pairings. Color contraction gives
$\sum_{A,B}(\delta^{AB})^2=N$, and therefore
\begin{equation}\label{eq:trace-free-scalar}
 \langle O(x)O(0)\rangle_0
 =\frac{2N}{16}\sum_{a,b,i,j}C_{abij}(x)^2
 =\frac{3N}{\pi^4|x|^8}.
\end{equation}
This computation applies to the actual reference table used in the
selected comparison. Indeed,
\eqref{eq:scribe-reference-dictionary} identifies its correlations with
Wick contractions for the potential kernel $D=(4\pi^2|x|^2)^{-1}$,
and \eqref{eq:scribe-reference-scalar-identification} identifies its
field $\mathcal O$ with precisely the ordered quarter-square field $O$
above. Thus the \emph{reference normalization} is a consequence of
that same-table dictionary and the calculation, rather than an
independent choice of a free comparison model:
\begin{equation}\label{eq:trace-reference-normalization}
 S_0(O,O;x_0,x_1)=\frac{3N}{\pi^4|x_0-x_1|^8}\qquad(x_0\neq x_1).
\end{equation}
The factor $1/4$ and the use of ordered
pairs are essential to the coefficient.

\subsection{Fixing the physical stress amplitude}
Here is the finite-dimensional step that prevents an arbitrary scalar
multiple of the stress tensor from entering the answer. All vector spaces
in this paragraph are real, including the real spaces of self-adjoint
tensor fields. Let $\mathcal K$ be the space of centered symmetric tensors
with the required rotation symmetry, let $\mathcal V$ be the space in
which their divergences lie, and let
$D:\mathcal K\to\mathcal V$ be the divergence map. The reference theory
has corresponding spaces $\mathcal K_0,\mathcal V_0$ and divergence $D_0$.
The following precise finite-dimensional data are used; they are the
properties under which Lemma~\ref{lem:ex-lattice-trace-rank-one} is
proved: $D_0$ is onto, $\dim\ker D_0=1$,
$\dim\mathcal K=\dim\mathcal K_0$,
and there is a surjective map
$Q:\mathcal K\to\mathcal V_0$ with $\ker Q=\ker D$.
Rank--nullity then gives
\begin{equation}\label{eq:trace-conserved-rank}
 \dim\ker D=\dim\mathcal K-\dim\mathcal V_0
           =\dim\mathcal K_0-\dim\mathcal V_0=1.
\end{equation}
In the construction the surjectivity of $Q$ is obtained from a surjective
finite-dimensional comparison composed with a test map into $\mathcal K$;
a surjective composition is sufficient to make $Q$ itself surjective.
The statements about $D_0$ are proved in
Proposition~\ref{prop:ex-trace-free-closure-rank}. The other properties are supplied by the conserved-line construction
of Section~\ref{sec:actual-stress-normalization}, on the completed
maps of the selected trace gallery, as follows. The map
$C_iB=\operatorname{center}(U_{i,4}B)$ is an isomorphism of the two
centered tensor spaces: its inverse centers $U_{i,4}^{-1}$, and the
unit corrections cancel because $U_{i,4}1=1$. This proves the dimension
identity. With $V_i$ the induced isomorphism of vector arrays,
$V_i^{-1}r_iDC_i$ converges to $D_0$ by the derivative-transport
argument there. If $R^{\rm right}$ is any linear right inverse of $D_0$, then
$V_i^{-1}r_iDC_iR^{\rm right}\to I$ and is invertible for large $i$. Thus
$Q_i=V_i^{-1}r_iD$ is onto and, since $r_i>0$ and $V_i$ is invertible,
$\ker Q_i=\ker D$. Choose any one such large $i$ for $Q$.
This construction uses only the uniform short-test estimate and
coherence of the gallery, not the trace asymptotic or its normalization.

Let $R\in\ker D$ be the physical centered stress tensor. At scales
$r_i>0$ let $U_i$ be the calibrated isomorphism from reference grade-four
fields to interacting grade-four fields. On tensors set
\[
 Y_i(T)_{ab}=r_i^4U_i^{-1}(T_{ab}).
\]
Suppose the comparison construction supplies eventually conserved tensors
$T_i\in\mathcal K$ and positive numbers $B_i$ with
\begin{equation}\label{eq:trace-trial-stress}
 Y_i(T_i)\longrightarrow R_0,\qquad
 B_i\operatorname{tr}Y_i(T_i)\longrightarrow4d_2O.
\end{equation}
In the selected construction these are the corrected tensors $T_i$ and
the arch scales $B_i=q_i$ of
Section~\ref{sec:ex-lattice-trace-scribe-producer}, with $R_0=M_0$ and
$4d_2=-\kappa$ there; the two limits are
\eqref{eq:ex-lattice-scribe-field-limit} and
\eqref{eq:ex-lattice-scribe-flux-limit}, which are conditional on the
inputs of that section.
The normalization test is a shell flux functional $\ell_i$ on tensors.
For example, fixing a separated spectator configuration $x$, a smooth
cutoff $\chi$ equal to one near that configuration, a direction $b$, and
one distinguished field $v$, its form is
\[
 \ell_i(T)=-r_i\sum_a\int(\partial_a\chi_i)(z)
   S_{T_{ab},v_i,\boldsymbol q_i}(z,r_ix)\,dz,
 \qquad \chi_i(z)=\chi(z/r_i).
\]
The derivatives of $\chi_i$ are supported away from the spectators, so the
integrals involve separated correlations. The translation Ward identity
for the physical tensor identifies this with a derivative insertion.
The needed convergence of that insertion and of the comparison shell
integral has the following exact consequence:
\begin{equation}\label{eq:trace-flux-test}
 \ell_i(R)\longrightarrow y\neq0,
 \qquad \ell_i(T_i)\longrightarrow y.
\end{equation}
The same nonzero $y$ occurs because $R_0$ has the reference translation
Ward identity. The detecting reference derivative field and the jet estimates
justifying these two limits are proved for the actual reference table
and the completed gallery charts in the proportionality argument of
Section~\ref{sec:ex-lattice-trace-scribe-producer}, where the two limits
appear as \eqref{eq:ex-lattice-scribe-common-test}.

\begin{lemma}\label{lem:trace-normalized-stress}
Equations \eqref{eq:trace-conserved-rank}--\eqref{eq:trace-flux-test} imply
\[
 Y_i(R)\longrightarrow R_0,\qquad
 B_i\operatorname{tr}Y_i(R)\longrightarrow4d_2O.
\]
\end{lemma}
\begin{proof}
The first limit in \eqref{eq:trace-flux-test} implies $R\neq0$.
Since the conserved space is one-dimensional, eventually $T_i=c_iR$
for a real scalar $c_i$. Linearity of $\ell_i$ gives
$c_i\ell_i(R)=\ell_i(T_i)$. For large $i$ the denominator $\ell_i(R)$
is nonzero, and the two limits give $c_i\to y/y=1$. In particular $c_i$
is eventually nonzero and $c_i^{-1}\to1$. Multiply both limits in
\eqref{eq:trace-trial-stress} by $c_i^{-1}$ and use linearity of $Y_i$
and of the trace. This yields exactly the asserted limits.
\end{proof}

Centering the trace in the chosen phase can differ from tracing the
mixed-state-centered tensor $R$ by a fixed multiple $c\mathbf1$ of the
unit. If $U_i\mathbf1_0=\mathbf1$ eventually and $B_ir_i^4\to0$
(for the selected data the latter is
\eqref{eq:ex-lattice-scribe-band-cost}), then
\[
 U_i^{-1}(B_ir_i^4c\mathbf1)=B_ir_i^4c\mathbf1_0\longrightarrow0.
\]
Consequently the selected centered trace $\Theta$ has the same limit
\begin{equation}\label{eq:trace-moving-label}
 a_i:=U_i^{-1}(B_ir_i^4\Theta)\longrightarrow4d_2O.
\end{equation}
The identities above explain how the physical Ward normalization, rather
than a later rescaling of $\Theta$, determines this coefficient.

\subsection{Rounded correlations and moving labels}
We next state exactly the remaining approximation data used to pass from
\eqref{eq:trace-moving-label} to the trace asymptotic. Write
$S_\chi(v,w;x_0,x_1)$ for the selected two-point function and
$S_0(v,w;x_0,x_1)$ for the reference one. Both are bilinear in the field
labels. The selected centered trace is a scalar under Euclidean rotations,
and translations preserve its correlations. Thus its two-point function
has the form
\begin{equation}\label{eq:trace-radial}
 S_\chi(\Theta,\Theta;x_0,x_1)=T(|x_0-x_1|)
       \quad(x_0\neq x_1).
\end{equation}
To justify this reduction, an orthogonal map sends the nonzero difference
$x_0-x_1$ to $-|x_0-x_1|e_0$; equal-length vectors are related by a
reflection, with the identity covering equality. A translation then sends
the two points to $-te_0/2,te_0/2$, where $t=|x_0-x_1|$.
Covariance and the scalar transformation law of $\Theta$ give
\eqref{eq:trace-radial}.

\begin{assumption}[Trace comparison along scales]\label{input:trace-scales}
For every sequence $s_n>0$ tending to zero there is a subsequence
$r_i=s_{n_i}$, with $n_i$ strictly increasing, and comparison data with
the following properties. The maps $U_i$ and numbers $B_i>0$ satisfy
\eqref{eq:trace-moving-label} and
\[
 \frac{\log(r_i^{-1})}{B_i}\longrightarrow\frac1\kappa,
 \qquad \kappa>0,\qquad d_2=-\frac\kappa4,
\]
where $\kappa$ is the slope of the selected running clock and the same
constant in its local response payment. Only this equality of the two
constants, not their separate numerical evaluation, enters the trace
coefficient.
There are rounded configurations $q_i(x)\to x$ and errors
$\varepsilon_i\to0$. Let $E_0$ be the reference grade-four field space
and fix its finite basis $(e_\alpha)$. At every coordinatewise separated
configuration $x$ in the interior of the comparison window, eventually
\begin{align}
 &\bigl|S_\chi(U_ie_\alpha,U_ie_\beta;r_iq_i(x))
                    -S_0(e_\alpha,e_\beta;q_i(x))\bigr|\nonumber\\
 &\hspace{30mm}\leq\varepsilon_i\bigl(1+M_{\alpha\beta}(q_i(x))\bigr),
 \label{eq:trace-basis-error}
\end{align}
where $M_{\alpha\beta}$ is continuous near $x$. The reference correlations
are continuous there. The comparison window includes a neighborhood of
$x(t)=(0,(t/2,t/2,t/2,t/2))$ for each $1/2<t<2$.
\end{assumption}
\begin{proposition}[The same-history scale supplier]
\label{prop:trace-scale-supplier}
Fix the common selected history and its extraction before choosing any
sequence of radii. Use its completed gallery, its physical stress and
centered trace, and the physical transfer
\eqref{eq:ex-lattice-scribe-physical-transfer} with the clock limit
\eqref{eq:ex-lattice-scribe-clock}. Then these objects supply
Input~\ref{input:trace-scales} for the same selected phase.
\end{proposition}
\begin{proof}
The scale estimates for the paid response, including the gain bound,
are supplied by Lemma~\ref{lem:trace-bill-scale-tickets}. The actual
coupling band is verified in
Section~\ref{sec:ex-lattice-trace-probe-assembly}. Its local atom
bounds are supplied by Subsection~\ref{subsec:trace-paid-atom-suppliers},
using the probe and cloud estimates of
Subsection~\ref{subsec:trace-eve-probe-cloud}; their thresholds are
included before the common history is chosen. Its grade-five buffer
uses the \hyperref[par:trace-actual-crown]{actual Crown witness}.
The conserved-line dimension is \eqref{eq:stress-physical-line},
with the same-table application explained above. Thus the local
estimates in the physical transfer are those of the selected objects,
not additional freely postulated responses or maps.

Let $s_n>0$ tend to zero. Apply
Lemma~\ref{lem:ex-lattice-firesky-selection} to this sequence. It
extracts strictly increasing indices $n_i$ and charts with radii
exactly $r_i=s_{n_i}$. The background history and its extraction
remain the already fixed ones; only charts and their indices are
selected after the sequence is given. All finite-grade requirements
use tails of these charts.

Take $U_i$ to be the completed grade-four map and $B_i=q_i$ to be
the actual summit arch scale. The physical transfer gives
$U_i^{-1}(B_i r_i^4\operatorname{tr}R)\to-\kappa O$.
The phase-centered trace differs from $\operatorname{tr}R$ by one
fixed scalar multiple of the unit. Unit preservation and
\eqref{eq:ex-lattice-scribe-band-cost} remove this difference after
pullback. Hence \eqref{eq:trace-moving-label} holds with
$4d_2=-\kappa$. The clock limit uses this same arch and same
positive slope $\kappa$, giving the required logarithmic ratio.
Neither the tensor nor its scalar normalization is reselected along
this subsequence.

Corollary~\ref{cor:trace-rounded-input} supplies the rounded
configurations, pairwise basis errors and common comparison window
for these maps and this phase. Its rounding bound tends to zero;
the reference kernels are smooth off the diagonals by their explicit
Wick formula. In particular they are continuous at every eligible
configuration, and the control functions there are bounded in a
neighborhood of that configuration. The window contains a
neighborhood of every $x(t)$ with $1/2<t<2$. These are all the
remaining clauses of Input~\ref{input:trace-scales}.
\end{proof}
The physical transfer in this proposition is the output of the
same-history response, sensor, branch and shell constructions of
Section~\ref{sec:ex-lattice-trace-scribe-producer}. The external
sensor and branch units are used there as distinct statements, with
their actual inputs matched separately. The proposition records the
order of choices needed at the endpoint; it does not replace any of
those quantitative constructions by compactness. The positive trace
spectral measure is constructed in Section~\ref{sec:trace-spectrum}
from this same phase and centered tensor.

\begin{lemma}\label{lem:trace-moving-correlation}
Assume Input~\ref{input:trace-scales}, the radial form
\eqref{eq:trace-radial} and the reference normalization
\eqref{eq:trace-reference-normalization}. Then, along each comparison
subsequence of the input and for each eligible fixed $x$,
\[
 r_i^8\log^2(r_i^{-1})T(r_i|q_i(x)_0-q_i(x)_1|)
       \longrightarrow \frac{A}{|x_0-x_1|^8},
 \qquad A=\frac{3N}{\pi^4}.
\]
\end{lemma}
\begin{proof}
The right side of \eqref{eq:trace-basis-error} tends to zero, since
$M_{\alpha\beta}(q_i(x))$ converges and is therefore bounded.
Continuity of $S_0$ gives convergence of the left correlation to
$S_0(e_\alpha,e_\beta;x)$. This also holds for convergent, varying labels:
expand each of them in the fixed finite basis. Every coordinate converges,
and bilinearity expresses the correlation as a finite sum of products of
two coordinates and one of the convergent basis correlations. In particular,
\eqref{eq:trace-moving-label} gives
\[
 S_\chi(U_ia_i,U_ia_i;r_iq_i(x))
 \longrightarrow (4d_2)^2 S_0(O,O;x).
\]
The definition of $a_i$, bilinearity and \eqref{eq:trace-radial} identify
the expression on the left with
$B_i^2r_i^8T(r_i|q_i(x)_0-q_i(x)_1|)$.
Here \eqref{eq:trace-radial} applies because the two points of
$r_iq_i(x)$ are distinct for large $i$: $q_i(x)\to x$ and the two points
of $x$ are distinct.
Multiply by $(\log(r_i^{-1})/B_i)^2$ and use
\eqref{eq:trace-reference-normalization}. The limiting coefficient is
\[
 \frac{(4d_2)^2}{\kappa^2}\frac{3N}{\pi^4}
       =\frac{(-\kappa)^2}{\kappa^2}\frac{3N}{\pi^4}
       =\frac{3N}{\pi^4}.
\]
All divisions are legitimate since $\kappa>0$ and $B_i>0$.
The common slope cancels before its numerical value is needed; this is
also the cancellation in \eqref{eq:scribe-slope-cancellation}.
\end{proof}

% Producer of the pointwise rounded comparison eq:trace-basis-error.
% Issue #336 (parent #56, C15). Correspondence and evidence:
% blueprint/units/YM-000026/handoff/results/C15/ADB991-claude-20261114-37f8/.
\subsection{The pointwise comparison at rounded configurations}
\label{subsec:trace-rounded-comparison}

We now derive the pointwise clause \eqref{eq:trace-basis-error} of
Input~\ref{input:trace-scales}, together with its window and continuity
clauses, from the selected gallery data of
Section~\ref{sec:ex-lattice-firesky-gallery}, the pairwise comparison
receipt contained in those data, and the phase construction of
Sections~\ref{ex-phases}--\ref{ex-phase-fields}. The derivation is
conditional on the inputs (G1)--(G4) and (Ph1)--(Ph3) listed below; the
only quantitative estimate it uses is the pairwise receipt (G3). For the selected funded family, (G3) is supplied by
Proposition~\ref{prop:trace-selected-rounded-receipt}. All other
steps are proved here or are results proved earlier in the text.

\paragraph{Configurations and the comparison window.}
For $n\ge1$ a configuration is $x=(x_0,\ldots,x_{n-1})\in(\mathbb
R^4)^n$, with the Euclidean norm of $\mathbb R^{4n}$. It is
\emph{coordinatewise separated} if for every coordinate direction
$a\in\{0,1,2,3\}$ the numbers $x_{0,a},\ldots,x_{n-1,a}$ are pairwise
distinct; such a configuration has distinct points. These configurations
form an open set. The comparison window is
\[
 \mathcal W=\{x\in(\mathbb R^4)^2:\ |x_0-x_1|\ge10^{-2},\ |x|\le100\}.
\]

\paragraph{Rounding.}
For $m>0$ and $x\in(\mathbb R^4)^n$ write $\lceil mx\rceil$ for the
coordinatewise ceiling, an element of $(\mathbb Z^4)^n$, and put
\[
 \rho_m(x)=m^{-1}\lceil mx\rceil .
\]
Each coordinate of $\rho_m(x)-x$ lies in $[0,m^{-1})$, so
\begin{equation}\label{eq:trace-rounding-error}
 |\rho_m(x)-x|\le \frac{2\sqrt n}{m}.
\end{equation}
Rounding is idempotent: $\lceil m\rho_m(x)\rceil=\lceil mx\rceil$,
because $m\rho_m(x)$ already has integer coordinates.

\paragraph{The selected gallery data.}
We use the following data, indexed by $i\in\mathbb N$, together with an
extraction order and an ultrafilter $\mathcal L$ on mesh--period pairs
$(N,p)$ finer than its fused filter, along which the mesh $N$ tends to
infinity.
\begin{enumerate}
\item[(G1)] \emph{Levels and radii.} A chart level $k_i$ with
 $k_i\to\infty$, a base arc $b_i>5$, and the radius
 \[
  r_i=a_{k_i}m_i,\qquad m_i=\mu_{k_i}^{\,b_i},
 \]
 where $a_k>0$ is the lattice spacing at depth $k$ and $\mu_k\ge1$ is the
 clock monitor, so that $r_i=g_{b_i}(k_i)$ in the notation of
 \eqref{eq:ex-lattice-firesky-gap}. The clock satisfies
 $\mu_k^5\to\infty$ as $k\to\infty$.
\item[(G2)] \emph{Chart readings and their limits.} For each grade $d$
 and each $i\ge d$, the grade-$d$ chart of level $i$ provides linear maps
 $X^{(i)}_{N,p}$ from the reference grade-$d$ field space $E_d$ into the
 raw grade-$d$ label space, one for every $(N,p)$. For each fixed $d$,
 eventually in $i$, these converge along $\mathcal L$ to a linear map
 $Y_{i}$, and the moving chart map of
 Input~\ref{input:trace-scales} at grade $d$ is
 $U_i=\mathsf P_dY_i$, where $\mathsf P_d$ is the quotient projection of
 the selected table. (This is the finite-dimensional extraction of
 Section~\ref{sec:ex-lattice-firesky-gallery}.)
\item[(G3)] \emph{Pairwise comparison receipt at cost
 $\varepsilon_i=(i+1)^{-1}$.} This is part of the one-step chart output
 recorded in Section~\ref{sec:ex-lattice-firesky-gallery}; its content is
 stated precisely in the next paragraph.
\item[(G4)] \emph{Chamber interiority.} There is a root depth
 $w_i\le k_i$ such that, if $x$ is fit at level $i$ (defined below), then
 for all sufficiently large meshes $N$ and every period $p$, every
 grade-$d$ insertion at a site $a_{k_i}\lceil m_ix\rceil_c$ belongs to the
 chamber-local domain of the centered tiling at depth $w_i$.
\end{enumerate}
A configuration $x\in(\mathbb R^4)^2$ is \emph{fit at level $i$} if the
grid configuration $a_{k_i}\lceil m_ix\rceil$ is coordinatewise
separated and each of its coordinates has absolute value at most one
eighth of the spacing of the middle cell of the chart. Fitness depends on
$x$ only through $\lceil m_ix\rceil$; by idempotence of rounding, $x$ is
fit if and only if $\rho_{m_i}(x)$ is fit.

\paragraph{The pairwise receipt.}
Fix $i$, a grade $d\le i$, a configuration $x$ fit at level $i$, and two
vectors $e,e'\in E_d$. Write $q=\rho_{m_i}(x)$. Put
\[
 j_0=S_0(e;q_0),\qquad j_1=S_0(e';q_1),\qquad
 \kappa=S_0(e,e';q_0,q_1)-j_0j_1,
\]
the reference one-point values and the reference connected two-point
value at the rounded points; these are defined because $q$ has distinct
points. For $(N,p)$ let $f_0^{N,p}$ and $f_1^{N,p}$ be the finite-volume
insertions of the labels $X^{(i)}_{N,p}e$ and $X^{(i)}_{N,p}e'$ at the
two lattice sites $a_{k_i}\lceil m_ix\rceil_0$ and
$a_{k_i}\lceil m_ix\rceil_1$. Let $\|\cdot\|_{\mathrm{root}}$ denote the
root test of the centered tiling at depth $w_i$ for the Wilson law at
$(N,p)$, as in
Section~\ref{sec:actual-phase-detector-moments}.
The receipt (G3) asserts: for every $i$, every $d\le i$, all basis vectors
$e,e'$ of $E_d$, all $N>k_i$ and all $p$, and every $x$ that is fit at
level $i$ with $q\in\mathcal W$,
\begin{align}
 \|f_c^{N,p}-j_c\|_{\mathrm{root}}&\le\varepsilon_i\qquad(c=0,1),
 \label{eq:trace-rounded-receipt-one}\\
 \|(f_0^{N,p}-j_0)(f_1^{N,p}-j_1)-\kappa\|_{\mathrm{root}}
 &\le\varepsilon_i .
 \label{eq:trace-rounded-receipt-two}
\end{align}
It also asserts the following \emph{width clause}: for every
$t\in(\mathbb Z^4)^2$ such that $y=m_i^{-1}t$ is coordinatewise separated
and $|y|\le100$, the configuration $y$ is fit at level $i$.

The receipt is formulated with a cumulant-corrected polynomial, and we
record why its three nonempty parts are exactly the left sides above. Let
$F$ be the reference moment profile at $q$, with
$F_\varnothing=1$, $F_{\{c\}}=j_c$ and $F_{\{0,1\}}=S_0(e,e';q_0,q_1)$.
The corrected polynomial $\Phi$ is the unique family
$(\Phi_s)_{s\subseteq\{0,1\}}$ with $\Phi_\varnothing=1$ and
\[
 \prod_{c\in s}f_c^{N,p}=\sum_{u\subseteq s}F_u\,\Phi_{s\setminus u}
 \qquad(s\subseteq\{0,1\}).
\]
Uniqueness and existence follow by induction on $|s|$, since the term
$u=\varnothing$ is $\Phi_s$. For $s=\{c\}$ the identity reads
$f_c^{N,p}=\Phi_{\{c\}}+j_c$. For $s=\{0,1\}$ it reads
\[
 f_0^{N,p}f_1^{N,p}=\Phi_{\{0,1\}}+j_0\Phi_{\{1\}}+j_1\Phi_{\{0\}}
     +S_0(e,e';q_0,q_1),
\]
and substituting the one-point parts gives
$\Phi_{\{0,1\}}=(f_0^{N,p}-j_0)(f_1^{N,p}-j_1)-\kappa$.

\paragraph{Phase inputs.}
Let $\mathcal A$, $\Omega$, $\Xi$ and $\mu$ be the commutative ground
algebra, the vacuum, the character space and the representing measure of
Sections~\ref{ex-phases} and~\ref{ex-phase-selection}, built from the
limiting (unselected) table of the same extraction order, and let
$S_\chi(w)=\chi(M(w))$ be as in \eqref{ex-phase-word}. We assume:
\begin{enumerate}
\item[(Ph1)] the hypotheses of that construction, in particular the
 separated-cluster limit of Lemma~\ref{ex-cluster-train}, the
 representation \eqref{ex-spectral-state} with full support, and the
 separating property of $\Omega$ for $\mathcal A$ used in the proof of
 Proposition~\ref{ex-character-measure};
\item[(Ph2)] the finite-volume hypotheses of the polynomial-packet
 argument of Section~\ref{sec:actual-phase-detector-moments}, from
 \eqref{eq:episode-cutoff-power} to \eqref{eq:episode-all-powers}:
 the planted inequality for the Wilson law at every $(N,p)$, the
 point-limit theorem for words with convergent labels along the selected
 order, and admission of the fixed reference depth; and
\item[(Ph3)] \emph{reality and null relations of the phases}: for every
 $\chi\in\Xi$ and every distinct-point word $w$,
 $S_\chi(\bar w)=\overline{S_\chi(w)}$, where $\bar w$ conjugates all
 labels without moving the points; and for the selected character $\chi$
 the phase correlations respect the null relations of the selected
 table, so that for raw labels $y,y'$ and distinct points $z_0,z_1$,
 \[
  S_\chi(\mathsf P_dy,\mathsf P_dy';z_0,z_1)=\chi\bigl(M(y@z_0\,y'@z_1)\bigr).
 \]
 (By Proposition~\ref{ex-no-new-null}, the null relations are respected
 under the detector hypotheses stated there.)
\end{enumerate}

\begin{lemma}[Conjugate words]\label{lem:trace-rounded-conjugate}
Under (Ph1) and the reality clause of (Ph3), $M(\bar w)=M(w)^*$ for every
distinct-point word $w$.
\end{lemma}
\begin{proof}
Put $Z=M(\bar w)-M(w)^*\in\mathcal A$. Every character is a star
homomorphism, so $\chi(Z)=S_\chi(\bar w)-\overline{S_\chi(w)}=0$ for
every $\chi\in\Xi$. By \eqref{ex-spectral-state} applied to $Z^*Z$,
\[
 \|Z\Omega\|^2=\langle\Omega,Z^*Z\Omega\rangle
   =\int_\Xi|\chi(Z)|^2\,d\mu(\chi)=0 .
\]
Since $\Omega$ is separating for $\mathcal A$, $Z=0$.
\end{proof}

Lemma~\ref{lem:trace-rounded-conjugate} supplies the identity
\eqref{eq:episode-conjugate}, which
Section~\ref{sec:actual-phase-detector-moments} takes as an input, in
the present setting; there it is used only to identify the limits of the
conjugated packets.

\begin{proposition}[Pointwise rounded comparison]
\label{prop:trace-rounded-comparison}
Assume (G1)--(G4) and (Ph1)--(Ph3), and let $\chi$ be the selected
character. Fix a grade $d$. There is $i_0$ such that for all $i\ge i_0$,
all basis vectors $e,e'$ of $E_d$, and every configuration $x$ that is fit
at level $i$ with $q=\rho_{m_i}(x)\in\mathcal W$,
\begin{equation}\label{eq:trace-rounded-comparison}
 \bigl|S_\chi(U_ie,U_ie';r_iq)-S_0(e,e';q)\bigr|
 \le\varepsilon_i\bigl(1+|S_0(e;q_0)|+|S_0(e';q_1)|\bigr).
\end{equation}
\end{proposition}
\begin{proof}
Choose $i_0\ge d$ such that the convergence in (G2) holds at grade $d$
for all $i\ge i_0$. Fix such $i$, the vectors $e,e'$ and $x$. Put
$y_0=Y_ie$, $y_1=Y_ie'$, so that
$X^{(i)}_{N,p}e\to y_0$ and $X^{(i)}_{N,p}e'\to y_1$ along
$\mathcal L$, and $\mathsf P_dy_0=U_ie$, $\mathsf P_dy_1=U_ie'$. Let $z$
be the grid configuration $a_{k_i}\lceil m_ix\rceil$. By (G1),
\[
 z=a_{k_i}m_i\,\rho_{m_i}(x)=r_iq .
\]
It is coordinatewise separated because $x$ is fit; in particular its
two points are distinct and lie on the depth-$k_i$ grid. Let $w$ be the
two-slot word carrying $y_0$ at $z_0$ and $y_1$ at $z_1$, and $w_c$ its
one-slot subwords.

\emph{Step 1: the two corrected packets.} Consider the polynomial
packets, in the sense of
Section~\ref{sec:actual-phase-detector-moments},
\begin{align*}
 &\text{words }[w_c],\varnothing\text{ with coefficients }1,-j_c
 &&(c=0,1),\\
 &\text{words }[w],[w_1],[w_0],\varnothing\text{ with coefficients }
 1,-j_0,-j_1,j_0j_1-\kappa .
\end{align*}
Their cutoff polynomials at $(N,p)$ are respectively $f_c^{N,p}-j_c$ and
\[
 f_0^{N,p}f_1^{N,p}-j_0f_1^{N,p}-j_1f_0^{N,p}+j_0j_1-\kappa
 =(f_0^{N,p}-j_0)(f_1^{N,p}-j_1)-\kappa ,
\]
and their compressed limits are
\[
 D_c=M(w_c)-j_cI,\qquad D_{01}=M(w)-j_0M(w_1)-j_1M(w_0)+(j_0j_1-\kappa)I .
\]
We check the hypotheses of the packet argument. The sites of every word
lie on the depth-$k_i$ grid and are coordinatewise separated, since $z$
is. The coefficients do not depend on $(N,p)$. The cutoff labels
$X^{(i)}_{N,p}e$ and $X^{(i)}_{N,p}e'$ converge along $\mathcal L$ to the
labels $y_0,y_1$ of the limiting words. The root depth satisfies
$w_i\le k_i$. Since $N\to\infty$ along $\mathcal L$, the set of $(N,p)$
with $N>k_i$ and $N$ beyond the threshold of (G4) belongs to
$\mathcal L$. On that set, (G4) puts $f_0^{N,p}$ and $f_1^{N,p}$ in the
chamber-local domain; constants, sums and products preserve that domain,
so all three cutoff polynomials belong to it. On the same set the receipt
\eqref{eq:trace-rounded-receipt-one}--\eqref{eq:trace-rounded-receipt-two}
bounds their root tests by $\varepsilon_i$. This is the root-price
hypothesis of the packet argument with $g=\varepsilon_i$ (and holds with
$g+\varepsilon$ for every $\varepsilon>0$ a fortiori).

By (Ph2) and Lemma~\ref{lem:trace-rounded-conjugate}, which supplies
\eqref{eq:episode-conjugate}, the packet argument applies to each of the
three packets and gives, by \eqref{eq:episode-all-powers},
\[
 \langle\Omega,(D^*D)^n\Omega\rangle\le\varepsilon_i^{2n}
 \qquad(n=0,1,2,\ldots;\ D=D_0,D_1,D_{01}).
\]

\emph{Step 2: every character.} Lemma~\ref{ex-character-moments}, which
uses the full-support representation of (Ph1), gives
$|\chi(D)|\le\varepsilon_i$ for each of the three elements and every
$\chi\in\Xi$. The two-point calculation following that lemma, with
$a=w_0$, $b=w_1$, $ab=w$, then gives
\[
 \bigl|S_\chi(w)-(\kappa+j_0j_1)\bigr|
 \le\varepsilon_i\bigl(1+|j_0|+|j_1|\bigr).
\]
For completeness: evaluating $\chi$ on $D_{01}$ and rearranging,
$S_\chi(w)-(\kappa+j_0j_1)=\chi(D_{01})+j_0\chi(D_1)+j_1\chi(D_0)$, and
the triangle inequality gives the bound.

\emph{Step 3: identification.} By the definition of $\kappa$,
$\kappa+j_0j_1=S_0(e,e';q_0,q_1)$. By the null-relation clause of (Ph3)
and Step 1,
\[
 S_\chi(w)=\chi(M(y_0@z_0\,y_1@z_1))
 =S_\chi(U_ie,U_ie';z_0,z_1)=S_\chi(U_ie,U_ie';r_iq).
\]
Substituting these identities and the definitions of $j_0,j_1$ proves
\eqref{eq:trace-rounded-comparison}. The threshold $i_0$ depends only on
$d$.
\end{proof}

\begin{corollary}[The pointwise clause of the trace input]
\label{cor:trace-rounded-input}
Assume (G1)--(G4) and (Ph1)--(Ph3), and that the selected chart
sequence realizes the comparison subsequence $r_i$ of
Input~\ref{input:trace-scales} at the grade $d=4$ of the reference
field space $E_0=E_4$. Assume also that the reference correlations are
continuous on configurations with distinct points. Put
$q_i(x)=\rho_{m_i}(x)$ and, for a basis $(e_\alpha)$ of $E_0$,
\[
 M_{\alpha\beta}(q)=|S_0(e_\alpha;q_0)|+|S_0(e_\beta;q_1)| .
\]
Then the pointwise clause of Input~\ref{input:trace-scales} holds with
these rounded configurations, with $\varepsilon_i=(i+1)^{-1}$, with this
$M_{\alpha\beta}$, and with comparison window $\mathcal W$: $q_i(x)\to x$
for every $x$, and at every coordinatewise separated $x$ in the interior
of $\mathcal W$, eventually in $i$, \eqref{eq:trace-basis-error} holds for
all $\alpha,\beta$. Moreover $M_{\alpha\beta}$ and the reference
correlations are continuous near such $x$, and the interior of
$\mathcal W$ contains a neighborhood of
$x(t)=(0,(t/2,t/2,t/2,t/2))$ for every $1/2<t<2$.
\end{corollary}
\begin{proof}
\emph{Convergence of the rounding.} Since $b_i>5$ and $\mu_k\ge1$,
$m_i=\mu_{k_i}^{b_i}\ge\mu_{k_i}^5$, which tends to infinity because
$k_i\to\infty$ (G1). Hence \eqref{eq:trace-rounding-error} gives
$|q_i(x)-x|\le2\sqrt n/m_i\to0$ for every configuration $x$ of $n$ points.

\emph{The estimate.} Let $x$ be coordinatewise separated and interior to
$\mathcal W$. The interior of $\mathcal W$ and the set of coordinatewise
separated configurations are open neighborhoods of $x$. Since
$q_i(x)\to x$, there is $i_1$ with $q_i(x)$ in both for $i\ge i_1$. For
such $i$ write $q_i(x)=m_i^{-1}t$ with $t=\lceil m_ix\rceil$. Then
$q_i(x)$ is coordinatewise separated and $|q_i(x)|\le100$, so the width
clause of (G3) shows that $q_i(x)$ is fit at level $i$; hence $x$ is fit,
by idempotence of rounding. Proposition~\ref{prop:trace-rounded-comparison}
applies for $i\ge\max(i_0,i_1)$ to every pair of basis vectors and gives
\eqref{eq:trace-basis-error} with the stated $\varepsilon_i$ and
$M_{\alpha\beta}$; note $\varepsilon_i\to0$.

\emph{Continuity.} The one-point functions $q\mapsto S_0(e_\alpha;q_0)$
and the two-point functions $q\mapsto S_0(e_\alpha,e_\beta;q)$ are
reference correlations of configurations with distinct points, hence
continuous on that open set by assumption; so is $M_{\alpha\beta}$, as a
sum of absolute values of continuous functions. The configuration $x$ has
distinct points, so these functions are continuous near $x$.

\emph{The window.} For $1/2<t<2$, the two points of $x(t)$ are at
distance $|x(t)_0-x(t)_1|=t>10^{-2}$, and $|x(t)|=t<100$. Both functions
$x\mapsto|x_0-x_1|$ and $x\mapsto|x|$ are continuous, so the strict
inequalities persist on a neighborhood of $x(t)$, which is therefore
contained in $\mathcal W$; it lies in the interior of $\mathcal W$.
Moreover $x(t)$ is coordinatewise separated, its coordinates being $0$
and $t/2\ne0$ in each direction.
\end{proof}

\paragraph{Status.}
Corollary~\ref{cor:trace-rounded-input} produces the pointwise,
continuity and window clauses of Input~\ref{input:trace-scales} for a
given comparison subsequence, conditional on (G1)--(G4) and
(Ph1)--(Ph3). The quantitative content is the pairwise receipt (G3) at
cost $(i+1)^{-1}$. Section~\ref{sec:ex-lattice-firesky-gallery} records
it as part of the one-step chart output and attributes the quantitative
estimates of that output to Section~\ref{sec:three-panel-weak-comparison}.
The root-test estimates themselves are supplied by
Proposition~\ref{prop:trace-selected-rounded-receipt} below: it attaches
the actual marked propagation and lawful representation at the root
depth to the level-budget estimate, and then purchases the cost
$(i+1)^{-1}$ and width clause by the selection thresholds.
The identification of the chart sequence with the comparison
subsequence of the input, the moving-label and logarithmic clauses, and
the reference normalization are separate inputs, discussed after
Input~\ref{input:trace-scales}.

% Producer of the pairwise receipt (G3) of subsec:trace-rounded-comparison.
% Issue #372 (parent #56, C15). Correspondence and evidence:
% blueprint/units/YM-000026/handoff/results/C15/author/hairer-claude-372-20261114/.
\subsection{The pairwise receipt from the level budget}
\label{subsec:trace-rounded-receipt}

We now derive the receipt (G3) of
Subsection~\ref{subsec:trace-rounded-comparison} from two properties of
the funded chart witnesses of
Section~\ref{sec:ex-lattice-firesky-gallery}: a scale allowance, stated in
(A) below, and a root-test bound at each level whose size is an explicit
negative power of the clock monitor, stated in (B) below. The width clause
uses only (A) and the growth of the monitor. The two root-test bounds
\eqref{eq:trace-rounded-receipt-one}--\eqref{eq:trace-rounded-receipt-two}
use only (B) and the growth of the monitor. The selection of the chart
levels then supplies (G3) at the cost $\varepsilon_i=(i+1)^{-1}$. At the
end of the subsection we show how the replica-root estimate of
Section~\ref{sec:rank-root-normalization} yields (B), and attach its inputs to the selected funded family.

\paragraph{Charts at a fixed level.}
Throughout, $\mu_k\ge1$ is the clock monitor of (G1), with
$\mu_k\to\infty$; it is the diverging monitor $\mathcal M_k$ of
Section~\ref{sec:rank-root-normalization}. Fix a grade cutoff $d$ and
let
\[
 M=M_d=d+2
\]
be the marker allowance of the funded chart family at grade $d$: the
maximal number of markers per cell in the replica estimate of
Section~\ref{sec:rank-root-normalization}. (It is not the arity cap $l$
of Section~\ref{sec:ex-lattice-firesky-gallery}.) Consider a chart of the
family at level $k$ with base arc $b>0$. Put
\[
 m=\mu_k^{\,b},\qquad r=a_km,
\]
so that $r=g_b(k)$ in the notation of \eqref{eq:ex-lattice-firesky-gap}.
Let $a_{\ell}$ be the spacing of the middle cell of the chart. Let $w$ be
the root depth of the chart: the depth of the centered tiling on which
the root tests of its comparison receipt are taken, and at which its
chamber condition is asserted. A configuration $x\in(\mathbb R^4)^2$ is
\emph{fit for the chart} if the grid configuration $a_k\lceil mx\rceil$
is coordinatewise separated and each of its coordinates has absolute
value at most $a_\ell/8$. For the chart selected at index $i$, with
$(k,b)=(k_i,b_i)$, this is exactly fitness at level $i$ as defined in
Subsection~\ref{subsec:trace-rounded-comparison}. As there, for a fit
$x$ and $e,e'\in E_{d'}$ with $d'\le d$, we write $f_0^{N,p}$,
$f_1^{N,p}$ for the finite-volume insertions of the chart labels of $e$
and $e'$ at the two sites of $a_k\lceil mx\rceil$. We write
$j_0,j_1,\kappa$ for the reference values at the rounded configuration
$\rho_m(x)$, and $\Phi_s$ ($s\subseteq\{0,1\}$) for the parts of the
corrected polynomial defined there. Thus $\Phi_{\{c\}}=f_c^{N,p}-j_c$ and
$\Phi_{\{0,1\}}=(f_0^{N,p}-j_0)(f_1^{N,p}-j_1)-\kappa$.

We use the following two properties of the funded chart witnesses. In
both, the level threshold is fixed before the chart is chosen.
\begin{enumerate}
\item[(A)] \emph{Allowance.} There is $K_A=K_A(d)$ such that every chart
 of the family at a level $k\ge K_A$ with base arc $b$ satisfies
 \begin{equation}\label{eq:trace-receipt-allowance}
  m\,\mu_k^{2}\le\frac{a_\ell}{a_k},
  \qquad\text{equivalently}\qquad r\,\mu_k^2\le a_\ell .
 \end{equation}
\item[(B)] \emph{Level budget.} There is $K_B=K_B(d)$ with the following
 property. Let $k\ge K_B$ and consider any chart of the family at level
 $k$, any grade $d'\le d$, any basis vectors $e,e'$ of $E_{d'}$, any
 $N>k$, any period $p$, and any $x$ fit for the chart with
 $\rho_m(x)\in\mathcal W$. Then, for every nonempty
 $s\subseteq\{0,1\}$, the root test of the centered tiling at the root
 depth $w$ of the chart, for the Wilson law at $(N,p)$, satisfies
 \begin{equation}\label{eq:trace-receipt-budget}
  \|\Phi_s\|_{\mathrm{root}}\le\mu_k^{-1/(16M)} .
 \end{equation}
\end{enumerate}
The thresholds $K_A$ and $K_B$ may depend on the grade cutoff, the chosen
bases and the window $\mathcal W$. They do not depend on the chart, on
$(N,p)$, or on $x$. Since there are finitely many grades $d'\le d$ and
finitely many pairs of basis vectors, a threshold for each fixed pair
yields one for all of them by taking the maximum. This chart-uniform form
is used in the proof of Proposition~\ref{prop:trace-rounded-receipt},
where the threshold must be known before the producer is called. A
weaker form, with a threshold allowed to depend on the level curve
returned by the producer, would suffice for the first assertion of that
proposition, but not for the argument given for its second assertion.

\begin{lemma}[The width clause]\label{lem:trace-receipt-width}
Let $R>0$. Consider a chart at a level $k\ge K_A$ with base arc $b$, and
suppose $\mu_k^2\ge8R$. Let $t\in(\mathbb Z^4)^2$ be such that
$y=m^{-1}t$ is coordinatewise separated and $|y|\le R$. Then $y$ is fit
for the chart. In particular, if $\mu_k^2\ge800$, the width clause of
(G3) holds for that chart.
\end{lemma}
\begin{proof}
Since $my=t$ has integer coordinates, $\lceil my\rceil=t$, and the grid
configuration of $y$ is
\[
 z=a_k\lceil my\rceil=a_kt=a_kmy=ry .
\]
Multiplication by $r>0$ preserves the strict order of real numbers, so
$z$ is coordinatewise separated because $y$ is. Each coordinate
satisfies $|z_{c,\nu}|=r|y_{c,\nu}|\le r|y|\le rR$, because a coordinate
of a vector in $\mathbb R^{8}$ is bounded by its Euclidean norm. By
\eqref{eq:trace-receipt-allowance}, which holds because $k\ge K_A$, and
by $\mu_k^2\ge8R$,
\[
 rR\le\frac{a_\ell}{\mu_k^2}\,R\le\frac{a_\ell}{8R}\,R=\frac{a_\ell}{8}.
\]
Hence $y$ is fit. The width clause of (G3) is the case $R=100$, and
there $8R=800$.
\end{proof}

\begin{lemma}[Cost of the root tests]\label{lem:trace-receipt-cost}
Assume (B), and let $0<\varepsilon\le1$. If $k\ge K_B$ and
$\mu_k\ge\varepsilon^{-16M}$, then for every chart of the family at level
$k$ the bounds
\eqref{eq:trace-rounded-receipt-one}--\eqref{eq:trace-rounded-receipt-two},
with the root test taken at the root depth $w$ of the chart, hold with
$\varepsilon$ in place of $\varepsilon_i$. This holds for every
grade $d'\le d$, all basis vectors $e,e'$ of $E_{d'}$, all $N>k$, all
$p$, and every $x$ fit for the chart with $\rho_m(x)\in\mathcal W$.
\end{lemma}
\begin{proof}
The function $t\mapsto t^{-1/(16M)}$ is decreasing on $(0,\infty)$, so
$\mu_k\ge\varepsilon^{-16M}$ gives
$\mu_k^{-1/(16M)}\le(\varepsilon^{-16M})^{-1/(16M)}=\varepsilon$. By
(B), $\|\Phi_s\|_{\mathrm{root}}\le\varepsilon$ for each nonempty
$s\subseteq\{0,1\}$. By the identities recalled above, the left sides
of \eqref{eq:trace-rounded-receipt-one} and
\eqref{eq:trace-rounded-receipt-two} are exactly
$\|\Phi_{\{c\}}\|_{\mathrm{root}}$ and
$\|\Phi_{\{0,1\}}\|_{\mathrm{root}}$.
\end{proof}

Since $\mu_k\to\infty$, the conditions
\[
 k\ge K_A(d),\qquad k\ge K_B(d),\qquad \mu_k^2\ge800,\qquad
 \mu_k\ge\varepsilon^{-16M_d}
\]
hold for all sufficiently large $k$. We write $K(d,\varepsilon)$ for the
least $K$ such that they hold for every $k\ge K$. (The monitor is not
assumed monotone, so this is the right form of the threshold.)

\begin{proposition}[Producer of the pairwise receipt]
\label{prop:trace-rounded-receipt}
Assume (A) and (B) at every grade cutoff. Identify the root depth $w_i$
of (G4) with the root depth $w$ of the chart selected at index $i$; that
is, (G4) is required at that depth. Suppose that for every index $i$ the
gallery selection of Section~\ref{sec:ex-lattice-firesky-gallery}, with
grade cutoff $d=i$ and cost index $l=i$, takes its chart at a level
\begin{equation}\label{eq:trace-receipt-level}
 k_i\ge K\bigl(i,(i+1)^{-1}\bigr).
\end{equation}
Then the selected data satisfy (G3) at the cost
$\varepsilon_i=(i+1)^{-1}$, including its width clause. Moreover
the selection can always be made so that \eqref{eq:trace-receipt-level}
holds for every $i$, without changing the conclusions of
Lemma~\ref{lem:ex-lattice-firesky-selection}.
\end{proposition}
\begin{proof}
Fix $i$, and let $M=M_i=i+2$. The selected chart has level
$k_i\ge K(i,(i+1)^{-1})$ and base arc $b_i$, so
$m=\mu_{k_i}^{b_i}=m_i$ and the fitness of the previous paragraphs is
fitness at level $i$. By the identification in the statement, the root
test of (G3) at index $i$ is taken at the depth $w_i=w$ of the selected
chart. Since $k_i\ge K_A(i)$ and $\mu_{k_i}^2\ge800$,
Lemma~\ref{lem:trace-receipt-width} with $R=100$ gives the width clause:
every $y=m_i^{-1}t$ that is coordinatewise separated with $|y|\le100$ is
fit at level $i$. Since $k_i\ge K_B(i)$ and
$\mu_{k_i}\ge(i+1)^{16M}=\varepsilon_i^{-16M}$,
Lemma~\ref{lem:trace-receipt-cost} with $\varepsilon=\varepsilon_i$ gives
\eqref{eq:trace-rounded-receipt-one}--\eqref{eq:trace-rounded-receipt-two}
at depth $w_i$. These hold for every grade $d'\le i$, all basis vectors
of $E_{d'}$, all $N>k_i$, all $p$, and every $x$ fit at level $i$ with
$\rho_{m_i}(x)\in\mathcal W$. As $i$ was arbitrary, this is (G3).

For the second assertion we use only the output of the one-step producer
as stated in Section~\ref{sec:ex-lattice-firesky-gallery}. The chart it
returns at level $k$ and base arc $b>5$ has exact radius equal to the
clock gap $g_b(k)=a_k\mu_k^b$ of \eqref{eq:ex-lattice-firesky-gap}, and the
producer accepts every threshold $\theta>0$ for that radius. Fix $i$ and
put $K=K(i,(i+1)^{-1})$; by (A) and (B) as stated, $K$ is determined
before the producer is called. Since $\mu_k\ge1$ and $b>0$, the radius
satisfies $g_b(k)\ge a_k$. Given the threshold $\theta>0$ required by the
selection at step $i$, call the producer instead with
\[
 \theta'=\min\Bigl(\theta,\ \min_{j<K}a_j\Bigr)>0
\]
(with $\theta'=\theta$ if $K=0$); this minimum is positive because
there are finitely many $j<K$ and each $a_j>0$. If the returned level $k$
satisfied $k<K$, the returned radius would satisfy
$g_b(k)\ge a_k\ge\theta'$, contradicting $g_b(k)<\theta'$. Hence
$k\ge K$, which is \eqref{eq:trace-receipt-level}. Since
$\theta'\le\theta$, the returned output also meets the threshold $\theta$
required in the proof of Lemma~\ref{lem:ex-lattice-firesky-selection},
and every other property of the output is unchanged. Thus that lemma
applies verbatim with the reduced thresholds, and its conclusions hold
for the resulting selection.
\end{proof}

\paragraph{Supply of the level budget.}
The bound \eqref{eq:trace-receipt-budget} is the replica-root estimate of
Section~\ref{sec:rank-root-normalization} at marker allowance $M$. Write
$q_0$ for the number of replica cells, which is also the site count of
the selected output torus. With
\[
 I_k=\frac{\log\mu_k}{16M},\qquad a=e^{2I_k}=\mu_k^{1/(8M)}\ge1 ,
\]
a coefficient whose pushed marked density has the local estimate with
coefficient price $a^{q_0|I|}$ and per-site expense $X\le1+M$ satisfies
\eqref{eq:rank-root}, whose right side is
$\exp\{P_R(X)-4D\log g/q_0\}/a^{|I|}$. The harmonic term
$-4D\log g/q_0$ is nonnegative and must be paid for. That section shows
\[
 P_R(X)-4D\log g/q_0\le T_M
\]
under the following conditions, all at the output level $m$ of the clock:
the coupling bound $\beta\le4\widehat\beta_m$; the harmonic condition
$8D\widehat\beta_m/q_0\le1$, which the clock supplies eventually; the
same clip and row bound $A$ as the clock stage, with
$\theta_R\le\theta_{\rm clock}$, $\Lambda\le\Lambda_{\rm clock}$ and
$\delta\le1$; and the late-scale purchases
$\widehat\beta_m\theta_{\rm clock}^3\le1$,
$\widehat\beta_m\theta_{\rm clock}^4\le1$ and
$\Lambda_{\rm clock}\theta_{\rm clock}^2\le1$. If moreover the stage ticket
satisfies $T_M\le I_k$, then the right side of \eqref{eq:rank-root} is
at most $e^{I_k}/a^{|I|}$. For a nonempty marker subset $|I|\ge1$ and
$a\ge1$, so this is at most $e^{I_k}/a=e^{-I_k}=\mu_k^{-1/(16M)}$. The
marker subsets $I$ are the nonempty $s\subseteq\{0,1\}$; since
$M\ge2$, they are admissible. The coefficients are the one-point and
corrected two-point parts $\Phi_s$ themselves, as
Section~\ref{sec:rank-root-normalization} requires.

The coefficient price $a^{q_0|I|}$ with $X\le1+M$ on the same pushed
density is Proposition~\ref{prop:allarity-output-price}, with $u=M$. It
is obtained from the scalar envelope \eqref{eq:allarity-density-envelope}
at the marker weight \eqref{eq:allarity-marker-weight}. That weight is
used from the launch at the middle level down to the output torus. The
leading interval above the middle level is run at a smaller weight. The
proposition's inputs, each on the common history of the chart witness,
are listed in the status paragraph at the end of
Section~\ref{sec:allarity-marked-propagation}.

\begin{proposition}[The receipt on the selected chart family]
\label{prop:trace-selected-rounded-receipt}
Use the common history and the funded panels of
Proposition~\ref{prop:trace-joint-chart-purchase}, with the ticket requests
\eqref{eq:rank-root-ticket-requests} included before the history is
chosen. For each fixed grade cutoff $d$, restrict to its three panels
and their offsets $t\in[0,1]$. Then (A) and (B) hold, with thresholds
uniform over those charts, all $N>k$, all periods, and all the fit
configurations in the stated window. The root test in (B) is at the
same depth $w$ as the chamber statement
Lemma~\ref{lem:ex-lattice-firesky-chamber}.
Consequently the selected gallery can satisfy (G3) at every index,
without changing its background history or extraction.
\end{proposition}
\begin{proof}
Fix $d$ and put $M=d+2$. Choose the constants of the leading
recurrence before its coupling threshold: after fixing its multiplier
and small-weight exponent $U$, take
$q\ge\max\{6,\kappa+4,\kappa+Un+2\}$, where $\kappa$ is the
upper step-count exponent and $n=2$ below. These are the initial,
spare-weight and launch margins of
Proposition~\ref{prop:allarity-launch}. The arcs of these panels have the fixed upper
bound $\bar b=\max_{0\le h\le2}(b_{d,h}+1)$. All applications below use
this bound; no uniformity over unbounded arcs is asserted. Put
$s=\widehat\beta_k$ and $H=\mu_k$. In the notation of
Section~\ref{sec:allarity-marked-propagation}, take $n=2$ marks per
replica cell and $u=M$. Thus $n\le u$, and its large marker weight is
$\omega=H^{1/(8M)}=a$. Its output level is the depth $w$ defined in
Subsection~\ref{subsec:ex-lattice-firesky-witness}, not the middle
level $\ell$. The output torus has exactly $q_0$ sites, one per cell
of this centered tiling.

\emph{The marked density and the five propagation inputs.}
We check the inputs listed at the end of
Section~\ref{sec:allarity-marked-propagation} on this particular
history. Use the corrected labels solving
\eqref{eq:ex-lattice-firesky-solves}, their actual insertions, and their
signed replica copies throughout.
\begin{enumerate}
\item The original mixed-grade marked expansion and its census payment
 in Section~\ref{sec:ex-lattice-firesky-gallery} give the terminal
 multi-marker representation of these insertions, including the
 ordinary and forced pieces. Its payment is
 \eqref{eq:trace-original-marked-payment}; the endpoint change of
 weights is \eqref{eq:trace-endpoint-marked-budget}. Both apply to
 the same supported terms, and retain the scalar and unmarked density.
 For the initial regular remainder, use the endpoint error estimate
 there with its parameter $g=q$ and its reserve order at least $q+4$:
 $H^{-D-8}A_0+d_s\le H^{-q}$. The forced payment requires a separate argument. Perform the
 endpoint regrouping before summing the marked factors, in the finite
 squarefree marker algebra with its weighted $\ell^1$ norm. This norm
 is submultiplicative, so the connected-set majorant and the gas
 bijection in Proposition~\ref{prop:restart-endpoint} apply unchanged.
 Every nonzero forced component contains a selected flag factor $X_z$:
 an old forced inventory contains an old flag, hence a new flag, and
 a mask-error factor vanishes unless its full support contains a new
 flag. Write $u_k=a_{\rm coast}\mathrm{hot}(s)^2$. The endpoint proof
 gives $|X_z|\le e^{-2u_k}$; reserve $e^{-u_k}$ in each component.
 The remaining site-factor price is precisely the $b_s$ of the
 gallery endpoint estimate. The old-support, marked-root and
 inherited-factor majorants are its $d_s,t_s,m_s$, evaluated in the
 same weighted norm. They do not have to be divided by $G_k$: their
 component already has the reserved flag factor. With the endpoint's
 two additional error terms their combined residual majorant is
 \[
  r_s=2H^{-3}+t_s+m_s+b_s\le H^{-1},\qquad
  Z(B)r_s\le1,\qquad e^{r_s}-1\le1.
 \]
 The connected sum is therefore at most
 $e^{-u_k}(e^{r_s}-1)\le e^{-u_k}=G_k$.
 This is a bound for the entire marked forced gas, not only its
 zero-marker part. At zero markers the original endpoint identity
 and scalar normalization are unchanged. The small weight
 $\omega_-\le1$ can only reduce the already paid marker norms.
 The terminal substitution of
 Proposition~\ref{prop:rank-selected-tested-readout} replaces only the
 near singleton by its masked raw polynomial and its two paid errors;
 the wide, forced and higher-mode terms are retained. It supplies the
 common polynomial and low-grade reading, rather than independently
 chosen maps. Its stamp and the calibrated polynomial price of
 Section~\ref{sec:allarity-coordinate-calibration} give the terminal
 source purchase and the principal power bound used by
 Lemma~\ref{lem:allarity-principal-power}.

 At every leading rung the actual compiler application is the one in
 the ``Room and the actual restarted steps'' argument of
 Section~\ref{sec:ex-lattice-firesky-gallery}. It keeps the path,
 clip, tube, forest, tag and positive loading. In particular
 Proposition~\ref{prop:restart-actual-closure} supplies the incoming
 neutral law and the genuine remainder bound; the joint-warning
 transfer is Lemma~\ref{lem:actual-neutral-warning}. These are the
 hypotheses of Section~\ref{sec:actual-local-compiler-handoff},
 not assumptions about an arbitrary coarse representation.

 For the local discrepancy use the guarded-law alternative
 Corollary~\ref{cor:allarity-guarded-purchase}, not exact transport
 by a normalized cutoff integral. The room and adjacent-profile
 comparison of Section~\ref{sec:ex-lattice-firesky-gallery} verify
 the hypotheses of Lemma~\ref{lem:allarity-pair-reader}. In particular
 the seed-free store, pose and enlarged guard rooms are those of
 Section~\ref{sec:actual-reader-dawn-selection}, transported with
 the actual paths, forest and tag. That lemma supplies one joint
 Gaussian reader array, its mean forecast and short covariance,
 with coordinate and covariance errors at most $H^{-f}$ and guard
 loss at most $H^{-2f}$, uniformly on the healthy coarse fields.

 Here is the component identification needed for this application.
 In the actual mixed compiler an entirely leading component has no
 body letter, and compatibility excludes background packets. Its
 integral is a product of the guarded principal values on the common
 Haar--Gaussian sample space. On a positive guard the field is
 healthy on the entire prescribed mask support, so that value is
 $p_I(\xi_{t,I})$. Whole-mode selection sums these products over
 partitions, with the same nonnegative forest/tag weights for every
 product, including the empty product, whose integral is one.
 These are exactly $M_S$ of Lemma~\ref{lem:allarity-guarded-gaussian}.
 Gauge invariance removes the patch gauge from each physical value.
 Finite Wick contraction with the same forecast and short covariance
 gives its $M^*_S$, the prescribed next polynomial moments.
 This identifies only the principal component, not the full
 interacting block density as Gaussian. The common loading is the
 one already fixed by the compiler; it is not replaced by a separate
 interacting partition function.

 The calibrated coefficient and covariance bounds give $A,B\le H^c$
 for a fixed $c$ on this finite menu. If $r$ is the degree of the
 polynomial in the guarded lemma, its connected error is
 $y\le C H^{cr-f}$. Fix the required remainder exponent $q$ first
 and choose $f>q+cr$. Then $2^nR_{e_*}H^qy\le1$ eventually.
 This is the local discrepancy purchase at the reserved support
 exponent. Realignment to the prescribed next polynomial propagates
 this error in the regular and forced remainders; no second history
 of normalized cutoff operators is iterated.
 Reader separation is Lemma~\ref{lem:allarity-halo-separation}.
\item The same reader comparison and the exact low-grade reading just
 identified give the moment error in
 Section~\ref{sec:allarity-cumulant-comparison}. The actual per-step
 compiler and terminal stamp give the healthy local error in
 Section~\ref{sec:marked-healthy-cumulant-launch}. The polynomial
 error estimates in Section~\ref{sec:ex-lattice-firesky-gallery}
 allow their fixed losses to be paid before the accuracy exponent is
 chosen, so both errors are at most $H^{-3/2}$ at the middle level.
 These are precisely the two hypotheses of
 Lemma~\ref{lem:allarity-corrected-pieces}; they concern the same
 principal polynomial and reference cumulant as in item 1.
\item Choose the support reserve before making the local compiler
 purchases: fix $0<\sigma\le1$, set $e_*=e+1+\sigma$ and $E=e+2$,
 and take $\alpha\ge8T_E$. Apply
 Proposition~\ref{prop:allarity-compiler-reserve} to the actual steps
 of item 1, with local support exponent $E+8$. The finite-stage
 block and entry purchases of
 Proposition~\ref{prop:trace-joint-chart-purchase} include these
 constants for each of its adjacent stage pairs. Thus the upper
 calculation and the terminal purchase hold at $e_*$ on the same
 letters. This verifies (E); Lemma~\ref{lem:allarity-spare-run} supplies the
 enlarged-exponent run, and
 Lemma~\ref{lem:allarity-remainder-restored} restores the marker
 weight. Coarse footprints contain the projected fine supports by
 the actual mixed compiler; Lemma~\ref{lem:allarity-reader-containment}
 gives the required containment. Halo separation supplies (S), so
 Proposition~\ref{prop:allarity-launch} applies.
\item The same reserved compiler supplies the lower calculation at
 exponent $e$ on $[w,\ell)$. Its remainder recurrence is
 Lemma~\ref{lem:allarity-lower-propagation}; its step count and
 forced inputs are supplied by Lemmas~\ref{lem:allarity-lower-count-free}
 and~\ref{lem:allarity-lower-forced}. The constant coupling, stage
 changes and normalizing scalar are those of
 Proposition~\ref{prop:restart-actual-closure}.
\item Corollary~\ref{cor:rank-trace-all-levels} supplies the lawful
 representations on that same course, at every level of $[w,k]$.
 At $w$ its clip and row bound are exactly those of $\pi_{J(w)}$,
 and it satisfies the coupling, threshold, loop and additive bounds
 required by Proposition~\ref{prop:rank-root-output-ticket}.
\end{enumerate}
These applications are uniform over the stated window: it is bounded
and its two points have distance at least $10^{-2}$. The finite-reader
and census estimates use these bounds, not a positive lower bound for
each coordinate difference. The latter differences need only be nonzero
for the chamber and ordering conditions. There are finitely many basis
pairs and grades through $d$. All the corresponding constants and all
three adjacent-stage cases are included in the finite-stage maxima
before the common history is chosen. The joint purchase then gives a
level tail independent of the chart offset, $N,p$ and $x$.

The propagated density is the push of
$\exp_*(-D_R)\prod_c(1+f_c^{N,p}z_c)$ times the actual background density,
with $D_R$ the reference cumulant polynomial. The field-independent
factor $\exp_*(-D_R)$ commutes with every push. The complete marked
push identities used in item 1 preserve every mode, not just the
singletons, and the unmarked coefficient preserves the partition
function. By Lemma~\ref{lem:allarity-invoice-parts} the three
nonempty coefficients are exactly the $\Phi_s$ of this subsection.
Proposition~\ref{prop:allarity-output-price} therefore gives, on this
lawful representation at $w$, the coefficient price $a^{q_0|s|}$ and
expense $X\le1+n\le1+M$.

\emph{Clock, ticket and uniform root depth.}
Lemma~\ref{lem:ex-lattice-firesky-clock-fields} supplies (A), the
plateau and the two-stage bracket uniformly for $b\le\bar b$.
The bracket gives $J(w)\ge J(k)-1$. In particular, for any fixed
level $m_0$, once $J(k)-1>J(m_0)$ every such chart has $w>m_0$.
This proves the uniform version of $w\to\infty$ needed for the
late-scale and harmonic purchases; one does not choose a different
tail for each chart curve. The ticket proof in
Proposition~\ref{prop:rank-root-output-ticket} is uniform in $w\le k$:
$J(w)\le J(k+1)$ lets the same countable payment buy
$T_M(\pi_{J(w)})\le I_k$. Combining it with the lawful representation
just obtained gives
\[
 \|\Phi_s\|_{\rm root}
 \le \frac{\exp\{P_{P_w}(X)-4D\log g/q_0\}}{a^{|s|}}
 \le \frac{e^{I_k}}{e^{2I_k|s|}}
 \le e^{-I_k}=\mu_k^{-1/(16M)}
 \qquad(s\ne\varnothing).
\]
This proves (B), including its finite-volume uniformity. Finally the
chamber lemma uses exactly the centered tiling at $w$. Its mesh
threshold is allowed in (G4); it does not weaken the root estimate,
which holds for every $N>k$. Proposition~\ref{prop:trace-rounded-receipt}
now supplies (G3), by decreasing the requested radius thresholds before
each chart selection. Only the selected charts change, not the history
or its extraction.
\end{proof}

\paragraph{Status.}
The conditional producer Proposition~\ref{prop:trace-rounded-receipt}
continues to apply to any family satisfying (A) and (B).
Proposition~\ref{prop:trace-selected-rounded-receipt} verifies those
properties on the selected funded family and identifies its root depth.
Thus (G3) is supplied for that family's rounded comparison. This
application retains the canonical selected-history and representation
constructions; it does not assert the receipt for an arbitrary chart
with merely the same numerical scales.



\subsection{Removing rounding and subsequences}
The rounded separation need not equal the intended radius. Pointwise
convergence at rounded arguments alone does not identify $T(r_i)$.
The positive spectral representation resolves this: $T$ is real,
nonnegative and decreasing, because $e^{-tp_0}$ has those properties for
each $p_0\geq0$. Integrability at the smaller positive time permits the
comparison under the integral.

\begin{theorem}\label{thm:trace-asymptotic}
Let $\Theta$ be the centered trace of Section~\ref{sec:trace-spectrum},
let $T$ be its two-point function in the selected phase as in
\eqref{eq:trace-radial}, and let $N=\dim_{\mathbb R}\mathfrak g$. Assume:
\begin{enumerate}
\item the reconstruction and stress hypotheses of
 Section~\ref{sec:trace-spectrum}, under which $T$ has the positive
 Laplace representation \eqref{eq:trace-positive-laplace};
\item invariance of the selected correlations under Euclidean translations
 and orthogonal transformations, together with the scalar law of
 $\Theta$, which give the radial form \eqref{eq:trace-radial};
\item the reference normalization
 \eqref{eq:trace-reference-normalization};
\item Input~\ref{input:trace-scales}.
\end{enumerate}
Then
\[
 \lim_{r\downarrow0}r^8\log^2(r^{-1})T(r)=\frac{3N}{\pi^4}>0.
\]
\end{theorem}
The reconstruction hypotheses and the selected-scale comparison are
matched to their producers below; the reference normalization has just
been proved for the actual comparator. The proof of the limiting passage
uses these statements separately.
\begin{proof}
First fix any one comparison subsequence $r_i$. Let
$c_i=r_i^8\log^2(r_i^{-1})\geq0$. For $1/2<t<2$ take the configuration
$x(t)$ in the input. Its Euclidean separation is exactly $t$, all four
coordinate separations are nonzero, and
$w_i(t):=|q_i(x(t))_0-q_i(x(t))_1|\to t$.
Lemma~\ref{lem:trace-moving-correlation} gives
\[
 c_iT(r_iw_i(t))\longrightarrow A/t^8.
\]
If $t>1$, eventually $w_i(t)>1$, so monotonicity yields
$c_iT(r_i)\geq c_iT(r_iw_i(t))$. If $t<1$, eventually
$0<w_i(t)<1$ and the reverse inequality holds. Given $\epsilon>0$,
choose $t_+>1$ and $t_-<1$ sufficiently close to one that
$|A/t_\pm^8-A|<\epsilon/2$. For all sufficiently large $i$, the two
comparison values are within $\epsilon/2$ of their limits. They bound
$c_iT(r_i)$ between $A-\epsilon$ and $A+\epsilon$. Thus
$c_iT(r_i)\to A$.

To obtain the limit at all positive radii, suppose it failed. There would
be $\epsilon>0$ and a sequence $s_n\in(0,1/n)$, $n\geq1$, for which
$|s_n^8\log^2(s_n^{-1})T(s_n)-A|\geq\epsilon$ for every $n$.
Input~\ref{input:trace-scales} supplies a comparison subsequence of this
same sequence. The preceding argument makes its weighted values converge
to $A$, contradicting that inequality. Finally $N>0$ for the nonabelian
simple Lie algebra under consideration, so $A>0$.
\end{proof}
Combining this theorem with Theorem~\ref{thm:profile-exclusion} proves the
finite-profile exclusion from the same trace comparison input. The application uses the same-history supplier
Proposition~\ref{prop:trace-scale-supplier} and the normalized stress
constructed above; the algebraic and limiting passages are separate
from those quantitative constructions.

% Issue #406 (parent #56, C15): the dilation current at short separation.
% Conditional as thm:trace-asymptotic, plus two stated assumptions; #337 open.
\begin{remark}[The dilation current at short separation]
\label{rem:trace-dilation-current}
Let $T_{ab}$, $\tau$ and $\Theta=\tau-c\mathbf1$ be as in
Section~\ref{sec:trace-spectrum}, let $T$ be the two-point function of
\eqref{eq:trace-radial}, and let $V$ be a continuously
differentiable vector field on $\mathbb R^4$ with
$\partial_aV_b+\partial_bV_a=2\varphi\delta_{ab}$ for a real function
$\varphi$. Assume in addition, in the selected phase and for
$x_0\neq x_1$, that $x_0\mapsto S_\chi(T_{ab},\Theta;x_0,x_1)$ is
continuously differentiable with
$\sum_a\partial_{(x_0)_a}S_\chi(T_{ab},\Theta;x_0,x_1)=0$ for each $b$
(the stress is conserved at separated points), and that
$S_\chi(\mathbf1,\Theta;x_0,x_1)=\langle\Theta\rangle_\chi$, as used in
Remark~\ref{rem:ex-selection-ground-only}. Section~\ref{sec:trace-spectrum}
does not assume the first; for the normalized stress of
Section~\ref{sec:physical-angular-actions} it follows, under their
hypotheses, from Propositions~\ref{prop:phase-ward-inheritance} (vary the
cutoff away from $x_1$) and~\ref{ex-phase-caps}. Put
$j_a(x_0)=\sum_bV_b(x_0)S_\chi(T_{ab},\Theta;x_0,x_1)$, the current
$\sum_bT_{ab}V_b$ at $x_0$ against the centered trace at $x_1$. Then
\[
 \sum_a\partial_{(x_0)_a}j_a(x_0)=\varphi(x_0)S_\chi(\tau,\Theta;x_0,x_1)
       =\varphi(x_0)T(|x_0-x_1|)\quad(x_0\neq x_1).
\]
Indeed, the product rule gives two sums. The one containing the
derivatives of $S_\chi(T_{ab},\Theta;x_0,x_1)$ vanishes by
conservation. The other is
$\sum_{a,b}\partial_aV_b\,S_\chi(T_{ab},\Theta;x_0,x_1)$; since
$T_{ab}=T_{ba}$, $\partial_aV_b$ may be replaced by
$\frac12(\partial_aV_b+\partial_bV_a)=\varphi\delta_{ab}$, and
bilinearity leaves $\varphi S_\chi(\tau,\Theta;x_0,x_1)$. Finally
$\tau=\Theta+c\mathbf1$, so bilinearity, the second assumption and
$\langle\Theta\rangle_\chi=0$ give
$S_\chi(\tau,\Theta;x_0,x_1)=S_\chi(\Theta,\Theta;x_0,x_1)$, which is
$T(|x_0-x_1|)$ by \eqref{eq:trace-radial}.

For the dilation field $V(z)=z$ one has $\varphi=1$, so for a pair at
distance $r$ the divergence of $j$ is $T(r)$. Under the hypotheses of
Theorem~\ref{thm:trace-asymptotic}, its conclusion
$r^8\log^2(r^{-1})T(r)\to3N/\pi^4>0$ as $r\downarrow0$, the limit required in
\eqref{ex-trace-limit}, gives $r_0>0$ with $T(r)>0$ for $0<r<r_0$.
Thus the divergence of $j$ is positive whenever
$0<|x_0-x_1|<r_0$: against the centered trace, the dilation current is
not divergence free at short separation. In fact $T(r)>0$ for every
$r>0$: in \eqref{eq:trace-positive-laplace} the measure is positive, and
nonzero since the limit is positive. Like the theorem, this
conclusion is conditional: it uses the four hypotheses of
Theorem~\ref{thm:trace-asymptotic} and the two assumptions above, none of
which is proved in this section. In the free Maxwell setting of
Section~\ref{sec:free-maxwell-radial-flux} the polarized stress of two
closed two-forms $f,g$ has zero trace, and its dilation current is
divergence free wherever $f$ and $g$ are Maxwell waves.
\end{remark}
